% !TeX root = ../navier-stokes-manuscript.tex
% ============================================================
% 2.1 Vortex Stretching
% ============================================================

The momentum equation has four pieces:
\[
  \underbrace{\partial_t u}_{\text{local acceleration}}
  +
  \underbrace{(u\cdot\nabla)u}_{\text{nonlinear transport}}
  =
  \underbrace{-\nabla p}_{\text{pressure}}
  +
  \underbrace{\nu\Delta u}_{\text{viscous diffusion}}.
\]
The central regularity problem lies in the competition between nonlinear transport and viscosity.
The hope is simple: viscosity smooths the fluid faster than nonlinear transport can sharpen it.
If that remains true at every scale, the velocity stays smooth.

\subsection{Vortex Stretching}\label{sub:ThePerfectlyStretchingVortex}

Along one material segment of a narrow vortex tube, the surrounding flow pulls in the axial direction. The symmetric part of the velocity gradient records that pull:
\[
  S:=\frac12\bigl(\nabla u+(\nabla u)^{\mathsf T}\bigr),
\]
and if \(e\) points along the segment while \(L(t)\) denotes its length, extension along \(e\) at rate \(\sigma(t)>0\) means
\[
  Se=\sigma(t)e,
  \qquad
  \sigma(t)>0,
  \qquad
  \frac{\dot L(t)}{L(t)}
  =
  e\cdot Se
  =
  \sigma(t).
\]
The same material piece is being drawn out along its own axis.

Because the fluid is incompressible, that extension cannot happen alone. If \(\mathcal V\) is the fixed material volume of the segment, \(A(t):=\mathcal V/L(t)\) its cross-sectional area, and \(r_{\mathrm{eff}}(t):=\sqrt{A(t)/\pi}\), then
\[
  \nabla\cdot u=0,
  \qquad
  \frac{d}{dt}(A L)=0,
  \qquad
  \frac{\dot A}{A}=-\sigma(t),
  \qquad
  \frac{\dot r_{\mathrm{eff}}}{r_{\mathrm{eff}}}
  =
  -\frac{\sigma(t)}{2}.
\]
Length is gained exactly where width is lost. The tube is becoming longer and thinner in one motion.

Its rotation moves inside that same deformation. With \(\omega:=\nabla\times u\), the vorticity obeys along the moving fluid
\[
  \frac{D\omega}{Dt}
  =
  S\omega+\nu\Delta\omega,
  \qquad
  \frac{D}{Dt}
  :=
  \partial_t+u\cdot\nabla,
\]
and when the vorticity aligns with the stretching direction, \(\omega=|\omega|e\),
\[
  S\omega
  =
  \sigma(t)\omega,
  \qquad
  \frac12\frac{D|\omega|^2}{Dt}
  =
  \sigma(t)|\omega|^2
  +
  \nu\,\omega\cdot\Delta\omega.
\]
Now the same positive strain that lengthens the segment also amplifies its aligned rotation, while viscosity enters the same balance without giving a pointwise sign by itself.

That growth sits inside the gradient even while the segment still carries only finite kinetic energy. If the speed stays bounded by \(U_0\) on the fixed-volume material segment \(\Omega_{\mathrm{seg}}(t)\), then
\[
  E_{\mathrm{seg}}(t)
  :=
  \frac12\int_{\Omega_{\mathrm{seg}}(t)}|u(x,t)|^2\,dx
  \leq
  \frac12U_0^2\mathcal V
  <
  \infty,
  \qquad
  \left|\frac12\bigl(\nabla u-(\nabla u)^{\mathsf T}\bigr)\right|_{\mathrm F}
  =
  \frac{|\omega|}{\sqrt2}
  \leq
  |\nabla u|_{\mathrm F}.
\]
So the same bounded-speed segment can keep finite energy while its rotating part strengthens and its radius falls. The gradients rise.

Fix \(t_0\) and let \(\ell:=r_{\mathrm{eff}}(t_0)\). The segment has narrowed to one width \(\ell\), and the unresolved question is whether passing to a finer width helps viscosity more than it helps the sharpening. For \(\lambda>1\), define
\[
  u_\lambda(x,t)
  :=
  \lambda u(\lambda x,\lambda^2t),
  \qquad
  p_\lambda(x,t)
  :=
  \lambda^2p(\lambda x,\lambda^2t).
\]
This rescaled field is another solution for comparison, not a later state of the same material tube. At time \(\lambda^{-2}t_0\), the corresponding vortex in \(u_\lambda\) has width \(\lambda^{-1}\ell\). Differentiating the rescaled fields gives
\[
\begin{aligned}
  \partial_tu_\lambda(x,t)
  &=
  \lambda^3(\partial_tu)(\lambda x,\lambda^2t),
  \\
  (u_\lambda\cdot\nabla)u_\lambda(x,t)
  &=
  \lambda^3\bigl((u\cdot\nabla)u\bigr)(\lambda x,\lambda^2t),
  \\
  \nabla p_\lambda(x,t)
  &=
  \lambda^3(\nabla p)(\lambda x,\lambda^2t),
  \\
  \nu\Delta u_\lambda(x,t)
  &=
  \lambda^3\nu(\Delta u)(\lambda x,\lambda^2t).
\end{aligned}
\]
At width \(\lambda^{-1}\ell\), nonlinear transport and viscosity are still locked together inside the equation: each has picked up the same \(\lambda^3\) factor, along with the other two terms. Narrowing the comparison has not given viscosity any relative advantage.

A finer comparison can still look larger or smaller simply because the ruler changes with the scale. Let \(\Lambda:=(-\Delta)^{1/2}\). At Sobolev height \(s\), a change of variables gives
\[
  \norm{u_\lambda(t)}_{\dot H^s}^2
  =
  \lambda^{2s-1}
  \norm{u(\lambda^2t)}_{\dot H^s}^2.
\]
The factor \(\lambda^{2s-1}\) comes from the rescaling itself, so it disappears only at \(s=\tfrac12\). Set
\[
  H(t)
  :=
  \frac12\norm{u(t)}_{\dot H^{1/2}}^2
  =
  \frac12\norm{\Lambda^{1/2}u(t)}_2^2.
\]
The levels on either side move in opposite directions:
\[
  \norm{u_\lambda(t)}_2^2
  =
  \lambda^{-1}\norm{u(\lambda^2t)}_2^2,
  \qquad
  H_\lambda(t)=H(\lambda^2t),
  \qquad
  \norm{\nabla u_\lambda(t)}_2^2
  =
  \lambda\norm{\nabla u(\lambda^2t)}_2^2.
\]
Kinetic energy falls under concentration, first-derivative energy rises, and this critical height does neither. Finite energy still leaves the concentration problem open. At \(H\), any increase has to come from the equation itself.

Follow that height along one solution. Apply \(\Lambda^{1/2}\) to the momentum equation and pair with \(\Lambda^{1/2}u\). Write the signed contribution of nonlinear transport as
\[
  \mathcal P_{1/2}(t)
  :=
  -\operatorname{Re}
  \pair
    {\Lambda^{1/2}u(t)}
    {\Lambda^{1/2}\bigl((u(t)\cdot\nabla)u(t)\bigr)}_{L^2}.
\]
The pressure pairing vanishes because \(\nabla\cdot u=0\), while the Laplacian contributes \(-\nu\norm{\Lambda^{3/2}u}_2^2\). What remains is
\[
  H'(t)
  +
  \nu\norm{\Lambda^{3/2}u(t)}_2^2
  =
  \mathcal P_{1/2}(t).
\]
The opening hope now has an exact form: \(H\) rises only when nonlinear production exceeds viscous dissipation. Under the same rescaling,
\[
  \mathcal P_{1/2}[u_\lambda](t)
  =
  \lambda^2\mathcal P_{1/2}[u](\lambda^2t),
  \qquad
  \nu\norm{\Lambda^{3/2}u_\lambda(t)}_2^2
  =
  \lambda^2\nu\norm{\Lambda^{3/2}u(\lambda^2t)}_2^2.
\]
The balance is still level at the critical height. What has changed is the clock: the same comparison is now being read at the contracted time coordinate \(t\mapsto\lambda^{-2}t\).

Viscosity itself carries that quadratic clock. For the linear heat equation \(v_t=\nu\Delta v\),
\[
  \widehat v(\xi,t+\delta t)
  =
  e^{-\nu|\xi|^2\delta t}\widehat v(\xi,t).
\]
A structure localized to spatial width \(r\) lives at frequencies \(|\xi|\simeq r^{-1}\), and the multiplier changes by order one when \(\nu|\xi|^2\delta t\simeq1\). That width is naturally paired with the viscous comparison time
\[
  \tau_\nu(r)
  :=
  \frac{r^2}{\nu}.
\]
This is not the lifetime of the material vortex. It is the time viscosity alone needs to alter a structure of width \(r\) by order one. At the finer comparison width,
\[
  \tau_\nu(\lambda^{-1}r)
  =
  \lambda^{-2}\tau_\nu(r),
\]
so the viscous comparison time contracts by the same factor as the time coordinate in the full \NavierStokes{} scaling. The unresolved spatial balance has survived again, but each narrower test now comes with less time.

Keep halving the width:
\[
  r_n:=2^{-n}r_0,
  \qquad
  \tau_n:=\frac{r_n^2}{\nu}.
\]
Then
\[
  \tau_n
  =
  \frac{r_0^2}{\nu}4^{-n},
  \qquad
  \sum_{n=0}^{\infty}\tau_n
  =
  \frac{4r_0^2}{3\nu}
  <
  \infty.
\]
So the comparison clocks attached to finer and finer widths fit inside a finite total time. That alone is still only a schedule of scales. For a singularity candidate, one actual solution would have to realize those widths,
\[
  r_0>r_1>r_2>\cdots\longrightarrow0
\]
at times
\[
  t_0<t_1<t_2<\cdots\uparrow T_*<\infty,
  \qquad
  t_{n+1}-t_n\asymp\tau_n,
\]
with comparison constants independent of \(n\). The remaining time would then satisfy
\[
  T_*-t_n
  \asymp
  \sum_{k=n}^{\infty}\tau_k
  =
  \frac{4r_n^2}{3\nu}.
\]
By stage \(n\), only order \(r_n^2/\nu\) of time is left in that same velocity history. The narrowing vortex has not yet produced a singularity, but it has converted the unresolved spatial competition into a temporal demand: if one smooth solution is really going to cross infinitely many shrinking widths before \(T_*\), the field must keep generating its next response on clocks that are collapsing just as fast.